% =============================================================================
% CompetitiveExam/CS301/09-efficient-sorting-answers.tex
% Answer key for 09-efficient-sorting-questions.tex.
%
% Real GATE PYQs Q1-Q6 are Tier-2 web-sourced and CoVe-checked inline (no
% forked claim-verifier exists in this environment). Independent mirrors:
%   Q1  GATE CSE 1996 Q40  — GeeksforGeeks + ExamSIDE worked answer.
%   Q2  GATE CSE 2014 S1 Q7 — GeeksforGeeks + ExamSIDE worked answer.
%   Q3  GATE CSE 2008 Q43  — GeeksforGeeks + GateQA (answer B).
%   Q4  GATE CSE 2019 Q32  — GeeksforGeeks (answer 0.08).
%   Q5  GATE CSE 2013 Q30  — official GATE 2013 CS paper Q.32 + pyqbox (C).
%   Q6  GATE IT 2008 Q41   — GeeksforGeeks + GATEOverflow (answer kn).
% Q7-Q10 are original; each was re-derived by hand this session (log below).
%
% Re-derivation log (2026-11-18):
%   Q1  first-element pivot; each partition compares pivot with all others:
%       sorted and reverse both remove one element per level -> n(n-1)/2. PASS
%   Q2  t1 = 4+3+2+1 = 10; t2 = 4 (pivot 4) + 2 (pivot 1) + 1 (pivot 3) = 7. PASS
%   Q3  worst split 1/5 : 4/5 -> T(n/5)+T(4n/5)+n; depth log_{5/4} n. PASS
%   Q4  worst slot = either extreme -> 2/25 = 0.08. PASS
%   Q5  k log k = Theta(log n) -> k = Theta(log n / log log n). PASS
%   Q6  base n: d = log_n(n^k) = k passes, each Theta(n) -> Theta(kn). PASS
%   Q7  [6,2,9,1,5,3] -> [2,6,9] + [1,3,5] -> [1,2,3,5,6,9]. PASS
%   Q8  median{3,2,1}=2; Lomuto -> [1,2,3,8,5,7], q=1, sizes 1 and 4. PASS
%   Q9  build -> [10,5,3,4,1,2]; extract -> [5,4,3,2,1]. PASS
%   Q10 16/8 = 2 passes; Theta(2(n+256)) time, Theta(n+256) space. PASS
%
% Compile with: cd CompetitiveExam/CS301 &&
%   TEXINPUTS=../../Preambles:$TEXINPUTS xelatex -interaction=nonstopmode 09-efficient-sorting-answers.tex
% =============================================================================

\documentclass[11pt]{article}
\usepackage[margin=1in]{geometry}
% =============================================================================
% Preambles/header.tex — shared LaTeX/Beamer preamble for this project
%
% Use from a Beamer lecture by prepending:
%     \input{header}
% and compiling with TEXINPUTS=../Preambles:$TEXINPUTS (the /compile-latex
% skill does this automatically).
%
% PALETTE CONTRACT. Color names defined here MUST match the names in
% Quarto/theme-template.scss so Beamer and Quarto renderings use the same
% palette. The scripts/check-palette-sync.sh script enforces this.
%
% To customize: edit the HEX values below AND the matching $variable in
% Quarto/theme-template.scss, then run scripts/check-palette-sync.sh.
% =============================================================================

% -----------------------------------------------------------------------------
% Engine detection + encoding
%
% Our /compile-latex skill uses XeLaTeX, where inputenc is unnecessary and
% can produce encoding warnings. Only load inputenc under pdfTeX.
% -----------------------------------------------------------------------------
\usepackage{iftex}
\ifPDFTeX
  \usepackage[utf8]{inputenc}
\fi

\usepackage{amsmath, amssymb, amsthm}
\usepackage{booktabs}
\usepackage{graphicx}

% -----------------------------------------------------------------------------
% PALETTE — mirrors Quarto/theme-template.scss variable names.
% Load xcolor and define colors BEFORE anything that references them
% (notably hyperref, hypersetup, and the Beamer color assignments).
% -----------------------------------------------------------------------------
\usepackage{xcolor}

\definecolor{primary-blue}{HTML}{012169}      % headings, accents
\definecolor{primary-gold}{HTML}{B9975B}      % emphasis, borders
\definecolor{highlight-yellow}{HTML}{F2A900}  % markers, alerts
\definecolor{light-bg}{HTML}{E8EDF5}          % title / section backgrounds
\definecolor{jet}{HTML}{1A1A1A}               % body text

% Semantic colors (used in diagrams and annotations)
\definecolor{positive}{HTML}{15803D}          % good / observed / identified
\definecolor{negative}{HTML}{B91C1C}          % bad / problematic / bias
\definecolor{neutral}{HTML}{525252}           % reference / context

% Highlight accents (match .hi-* classes in the SCSS)
\definecolor{hi-slate}{HTML}{314F4F}
\definecolor{hi-green}{HTML}{2E8B57}
\definecolor{hi-red}{HTML}{C41E3A}

% Semantic aliases so skill templates and snippets can write
% \textcolor{accent}{...} without hard-coding "primary-blue".
\colorlet{accent}{primary-blue}
\colorlet{accent2}{primary-gold}

% -----------------------------------------------------------------------------
% Hyperref — loaded AFTER xcolor + palette so link colors resolve.
% -----------------------------------------------------------------------------
\usepackage{hyperref}
\hypersetup{colorlinks=true, linkcolor=primary-blue, urlcolor=primary-blue,
            citecolor=primary-blue, pdfborder={0 0 0}}

% -----------------------------------------------------------------------------
% Course code — set once per deck (after \input{header}, before \begin{document})
% via \coursecode{CS401}. Shown in the footer on every slide so a deck's course
% is identifiable without opening the title page. Empty by default.
% -----------------------------------------------------------------------------
\makeatletter
\newcommand{\coursecode}[1]{\gdef\@coursecodetext{#1}}
\gdef\@coursecodetext{}
\makeatother

% -----------------------------------------------------------------------------
% Beamer theme assignments (only apply under Beamer).
%
% We detect Beamer via \@ifclassloaded{beamer}, which is the documented,
% non-internal way. This must be wrapped in \makeatletter/\makeatother
% because \@ifclassloaded contains an @-catcode character.
% -----------------------------------------------------------------------------
\makeatletter
\@ifclassloaded{beamer}{%
  \usefonttheme{professionalfonts}
  \setbeamercolor{structure}{fg=primary-blue}
  \setbeamercolor{frametitle}{fg=primary-blue}
  \setbeamercolor{title}{fg=primary-blue}
  \setbeamercolor{subtitle}{fg=primary-gold}
  \setbeamercolor{author}{fg=primary-blue}
  \setbeamercolor{institute}{fg=neutral}
  \setbeamercolor{date}{fg=primary-gold}
  \setbeamercolor{itemize item}{fg=primary-blue}
  \setbeamercolor{itemize subitem}{fg=primary-gold}
  \setbeamercolor{alerted text}{fg=primary-gold}
  \setbeamercolor{block title}{fg=white, bg=primary-blue}
  \setbeamercolor{block body}{bg=light-bg}
  % Semantic block variants: block=definition/concept (blue), exampleblock=
  % worked example (green/positive), alertblock=key takeaway or caution (gold).
  % Keep to <=2 colored boxes per slide across ALL three variants (INV-7).
  \setbeamercolor{block title example}{fg=white, bg=positive}
  \setbeamercolor{block body example}{bg=light-bg}
  \setbeamercolor{block title alerted}{fg=white, bg=primary-gold}
  \setbeamercolor{block body alerted}{bg=light-bg}

  % Minimal footer: course code (left) + page X / N (right), no navigation clutter
  \setbeamertemplate{navigation symbols}{}
  \setbeamertemplate{footline}{%
    \color{neutral}\footnotesize\hspace{0.7em}\@coursecodetext\hfill
    \insertframenumber/\inserttotalframenumber\hspace{0.7em}\vspace{0.3em}}
}{}
\makeatother

% -----------------------------------------------------------------------------
% TikZ setup — libraries the snippet gallery uses
% -----------------------------------------------------------------------------
\usepackage{tikz}
\usetikzlibrary{arrows.meta, positioning, calc, decorations.pathreplacing,
                fit, shapes.geometric, backgrounds}

% Shared TikZ styles — snippets and hand-written diagrams can pick these up
% via `\begin{tikzpicture}[every node/.style={...}]` or by naming specific
% styles (e.g., `\node[dag-node] (X) at (0,0) {X};`).
\tikzset{
  % Node styles for causal DAGs and similar
  dag-node/.style = {
    draw=primary-blue, fill=light-bg, rounded corners=2pt,
    minimum width=1.2cm, minimum height=0.9cm, align=center, font=\small
  },
  decision-node/.style = {
    draw=primary-gold, fill=white, diamond, aspect=1.8,
    minimum width=1.8cm, minimum height=1.2cm, align=center, font=\small,
    inner sep=1pt
  },
  % Edge styles — observed vs counterfactual
  observed-edge/.style = {-{Latex[length=2.5mm]}, thick, draw=jet},
  counterfactual-edge/.style = {-{Latex[length=2.5mm]}, thick, draw=neutral, dashed},
  confound-edge/.style = {-{Latex[length=2.5mm]}, thick, draw=negative, dashed},
  % Data-point styles for plots
  observed-dot/.style = {fill=primary-blue, draw=primary-blue, circle, inner sep=1.2pt},
  counterfactual-dot/.style = {fill=white, draw=neutral, circle, inner sep=1.2pt}
}

% -----------------------------------------------------------------------------
% Convenience macros
% -----------------------------------------------------------------------------
\newcommand{\muted}[1]{\textcolor{neutral}{#1}}
\newcommand{\key}[1]{\textcolor{primary-gold}{\textbf{#1}}}
\newcommand{\good}[1]{\textcolor{positive}{#1}}
\newcommand{\bad}[1]{\textcolor{negative}{#1}}

% Standout frame macro: full-bleed dark background for section transitions.
% Beamer-only (uses \begin{frame}/\setbeamercolor/\usebeamerfont) -- do not
% call from an article-class file (e.g. Notes/<CODE>/*-notes.tex). \muted,
% \key, \good, \bad, and \coursecode above are plain xcolor/gdef and are
% safe to use from either class; verified when Notes/article-class support
% was added.
% Expand: \transitionslide{Your transition title}
\newcommand{\transitionslide}[1]{%
  \begingroup
  \setbeamercolor{background canvas}{bg=primary-blue}
  \begin{frame}[plain]
    \centering
    \vfill
    {\usebeamerfont{title}\color{white}\Large #1\par}
    \vfill
  \end{frame}
  \endgroup
}

% Section divider frame (Beamer-only), an alternative to \transitionslide with a
% darker two-line style: near-black (jet) background, a small white label line
% above a larger gold title, and a thin gold rule along the bottom edge.
% Expand: \sectiondivider{Section 2}{Pointers}
\newcommand{\sectiondivider}[2]{%
  \begingroup
  \setbeamercolor{background canvas}{bg=jet}
  \begin{frame}[plain]
    \vspace*{3.0cm}
    {\color{white}\ttfamily\large #1\par}
    \vspace{0.5cm}
    {\color{primary-gold}\LARGE #2\par}
    \begin{tikzpicture}[remember picture, overlay]
      \draw[primary-gold, line width=2.5pt]
        ($(current page.south west)+(0,0.1)$) -- ($(current page.south east)+(0,0.1)$);
    \end{tikzpicture}
  \end{frame}
  \endgroup
}

% -----------------------------------------------------------------------------
% Channel watermark — "Concepts That Click" (YouTube).
%
% Article-class files (CompetitiveExam Q/A sets, Notes, Assignments) get a
% light diagonal draft watermark on every page plus a centered footer naming
% the channel. Beamer decks get a subtle TikZ background watermark instead
% (the footline already carries the course code). Centralized here so every
% MCQ PDF is branded without per-file edits.
% -----------------------------------------------------------------------------
\makeatletter
\@ifclassloaded{article}{%
  \usepackage{draftwatermark}
  \SetWatermarkText{Concepts That Click}
  \SetWatermarkScale{1}
  \SetWatermarkFontSize{2cm}
  \SetWatermarkLightness{0.86}
  \SetWatermarkAngle{45}
  \usepackage{fancyhdr}
  \pagestyle{fancy}
  \fancyhf{}
  \fancyfoot[C]{\footnotesize\textcolor{neutral}{Concepts That Click~|~YouTube: \href{https://www.youtube.com/@concepts.thatclick}{@concepts.thatclick}}}
  \renewcommand{\headrulewidth}{0pt}
  \renewcommand{\footrulewidth}{0pt}
  \fancypagestyle{plain}{%
    \fancyhf{}%
    \fancyfoot[C]{\footnotesize\textcolor{neutral}{Concepts That Click~|~YouTube: \href{https://www.youtube.com/@concepts.thatclick}{@concepts.thatclick}}}%
    \renewcommand{\headrulewidth}{0pt}%
    \renewcommand{\footrulewidth}{0pt}%
  }
}{}
\@ifclassloaded{beamer}{%
  \setbeamertemplate{background}{%
    \begin{tikzpicture}[remember picture, overlay]
      \node[rotate=45, opacity=0.055, align=center]
        at (current page.center)
        {\Huge\textcolor{neutral}{Concepts That Click}};
    \end{tikzpicture}%
  }
}{}
\makeatother 
\coursecode{CS301}
\renewcommand{\thesection}{9.\arabic{section}}
\setlength{\emergencystretch}{2.5em}

\title{Competitive Exam Practice 9: Efficient Sorting\\\large Answer Key}
\author{Auqib Hamid Lone\\[0.3em]
  \emph{Department of Computer Science \& Engineering}, \emph{GCET Ganderbal}}
\date{\today}

\begin{document}
\maketitle

\noindent Answers with full working. MCQ option letters refer to the student question paper.

\section*{Q1 --- Answer (C)}

With the first element as pivot, partitioning a subarray of $m$ elements compares the pivot with each of the other $m-1$ elements. For the sorted input, every pivot is the \emph{smallest} element of its subarray, so the left part is empty and the right part has $m-1$ elements. For the reverse input, the first pivot $n$ is the largest, the next pivot $n-1$ is the largest remaining, and every subsequent pivot is an extreme; each partition again removes only the pivot. In both runs the subarray sizes are $n,n-1,\dots,1$, so
\[
  C_1 = C_2 = (n-1)+(n-2)+\cdots+1 = \frac{n(n-1)}{2}.
\]
The equal count is why this input is the worst case, not a difference between the two orders.

\section*{Q2 --- Answer (C)}

For $\{1,2,3,4,5\}$ every pivot is the minimum: partitions cost $4,3,2,1$ comparisons, so $t_1 = 4+3+2+1 = 10$. For $\{4,1,5,3,2\}$, pivot $4$ is compared with all four others (4 comparisons) and leaves $\{1,3,2\}$ and $\{5\}$; the subarray $\{1,3,2\}$ is partitioned around $1$ with $2$ comparisons, leaving $\{3,2\}$, which takes $1$ more. Hence $t_2 = 4+2+1 = 7$. Therefore $t_1 > t_2$, option (C).

\section*{Q3 --- Answer (B)}

The guaranteed split puts at least $1/5$ of the elements on each side, so in the worst case one side is $n/5$ and the other is $4n/5$. Partitioning costs at most $n$ comparisons, giving
\[
  T(n) \le T(n/5) + T(4n/5) + n.
\]
Unrolling, the depth is $\log_{5/4} n$, and each level costs at most $n$, so $T(n)=O(n\log n)$. Option (A) treats both sides as $n/5$ (loses $3/5$ of the data), (C) treats both as $4n/5$ (over-counts), and (D) assumes a balanced split that the procedure does not promise.

\section*{Q4 --- Answer 0.08}

In a partition of $25$ distinct elements, the pivot is in its final position. The worst possible location is either extreme (first or last), because that leaves a subproblem of size $24$ and one of size $0$. The pivot is chosen uniformly from $25$ positions, so the probability is $2/25 = 0.08$.

\section*{Q5 --- Answer (C)}

Heap sort on $k$ elements costs $\Theta(k\log k)$ in the worst case. Setting $k\log k = \Theta(\log n)$ gives $k = \Theta(\log n/\log\log n)$: with $k$ that size, $k\log k \approx (\log n/\log\log n)(\log\log n) = \log n$. Option (D) ignores the $\log k$ factor, (A) is far too small, and (B) is not the solution of the equation.

\section*{Q6 --- Answer (B)}

Take base $b=n$. A number at most $n^k$ has $\lceil\log_n(n^k)\rceil = k$ base-$n$ digits, so there are $k$ passes. Each pass is a stable counting sort over $n$ elements and $n$ buckets, costing $\Theta(n+n)=\Theta(n)$. The total is $\Theta(kn)$. (With a fixed base the digit count carries a $\log n$ factor; the option list corresponds to the base-$n$ analysis.)

\section*{Q7 --- Split tree and merges}

Split: $[6,2,9,1,5,3]\to[6,2,9]\;[1,5,3]\to[6,2]\;[9]\;[1,5]\;[3]\to[6]\;[2]\;[9]\;[1]\;[5]\;[3]$.

Merge level 1: $[6]+[2]\to[2,6]$; then $[2,6]+[9]\to[2,6,9]$. On the right, $[1]+[5]\to[1,5]$; then $[1,5]+[3]\to[1,3,5]$.

Merge level 2: $[2,6,9]+[1,3,5]\to[1,2,3,5,6,9]$.

There are $\lceil\log_2 6\rceil = 3$ levels of merging. A sorted input does \emph{not} reduce the comparison count: merge sort still walks both runs and copies every element once per level, so the cost stays $\Theta(n\log n)$.

\section*{Q8 --- Median-of-three partition}

The three candidates are $A[0]=3$, $A[2]=2$, $A[5]=1$; their median is $2$. Swapping the median to the last position gives $[3,7,1,8,5,2]$, so the pivot is $2$.

Lomuto scan with $i=-1$: $3>2$ (no move); $7>2$ (no move); $1\le2$, so $i=0$ and swap $A[0],A[2]$: $[1,7,3,8,5,2]$; $8>2$ and $5>2$ (no moves). Final swap $A[1],A[5]$: $[1,2,3,8,5,7]$, returning index $1$. The left side $[1]$ has size $1$; the right side $[3,8,5,7]$ has size $4$.

\section*{Q9 --- Build and extract}

Build from $[4,10,3,5,1,2]$; the last parent is index $2$. Sift index $2$ ($3$ with child $2$): no change. Sift index $1$ ($10$ with children $5,1$): no change. Sift index $0$ ($4$ with children $10,3$): swap with $10$, giving $[10,4,3,5,1,2]$; continue at index $1$: swap $4$ with $5$, giving $[10,5,3,4,1,2]$; index $3$ is a leaf. The heap is $[10,5,3,4,1,2]$.

Extract the maximum: swap $10$ with the last element $2$ and shrink the heap to $[2,5,3,4,1]$; sift down index $0$: swap $2$ with $5$ $\to[5,2,3,4,1]$; sift index $1$: swap $2$ with $4$ $\to[5,4,3,2,1]$; index $3$ is a leaf. The result $[5,4,3,2,1]$ satisfies the max-heap property ($5\ge4,3$; $4\ge2,1$).

\section*{Q10 --- Radix design}

(a) With $8$-bit digits, $16/8 = 2$ passes.

(b) Radix sort is $\Theta(d(n+k))$ time and $\Theta(n+k)$ space for $d$ digits and base $k$. Here $d=2$, $k=256$, so time is $\Theta(2(n+256))$ and auxiliary space is $\Theta(n+256)$.

(c) The $\Omega(n\log n)$ bound applies only to algorithms that learn about keys by comparing them. Radix sort reads each digit and never compares two whole keys, so the decision-tree argument --- and hence the bound --- does not apply.

\end{document}
